% Generated by tex/build_swebok_structure.py from tex/chapters/ch01/05_要求分析.tex
\subsubsection{基本的な要求分析}

よい要求には、個々の要求として望ましい性質と、要求集合として望ましい性質がある。

\paragraph{個々の要求の望ましい性質}

\par\smallskip\noindent\refstepcounter{table}\label{tab:ch01-08}\textbf{表 \thetable: 第01章 ソフトウェア要求 表08}\par\smallskip
\begin{tabularx}{0.97\linewidth}{@{}p{.26\linewidth}X@{}}
\toprule
性質 & 説明 \\
\midrule
\term{一意に解釈できる} & 読む人によって意味が変わらない。 \\
\term{テスト可能である} & 満たしたかどうかを確認できる。必要なら数値で示す。 \\
\term{拘束力がある} & 顧客がそれに価値を認め、満たされないことを受け入れない。 \\
\term{原子的である} & 一つの要求が一つの判断だけを表す。 \\
\term{真のニーズを表す} & 単なる思いつきや特定解ではなく、実際の問題に基づく。 \\
\term{利害関係者の語彙を使う} & 利用者や業務専門家が理解できる言葉で書く。 \\
\term{関係者に受け入れられる} & 重要な利害関係者が納得できる。 \\
\bottomrule
\end{tabularx}

\paragraph{要求集合の望ましい性質}

\par\smallskip\noindent\refstepcounter{table}\label{tab:ch01-09}\textbf{表 \thetable: 第01章 ソフトウェア要求 表09}\par\smallskip
\begin{tabularx}{0.97\linewidth}{@{}p{.26\linewidth}X@{}}
\toprule
性質 & 説明 \\
\midrule
\term{完全性} & 通常ケースだけでなく、境界条件、例外条件、セキュリティ上の必要事項も扱う。 \\
\term{簡潔性} & 不要な記述や重複が少ない。 \\
\term{内部整合性} & 要求同士が矛盾しない。 \\
\term{外部整合性} & 契約、法令、標準、業務文書などの外部資料と矛盾しない。 \\
\term{実現可能性} & 費用、納期、人員、技術制約の範囲で実現できる。 \\
\bottomrule
\end{tabularx}

\MiniBox{5つのなぜ}{要求が解決策の形で出てきたときは、「なぜそれが必要なのか」を繰り返し問う。目的は、相手が本当に解決したい問題に到達することである。}
